\chapter{Trigger Point Injection}

\section*{Overview}
Trigger point injection is a procedure in which a local anesthetic, often with a corticosteroid, is injected directly into a myofascial trigger point. It is used to relieve pain, muscle spasm, and referred tenderness associated with myofascial pain syndrome.

\section*{Alternative Names}
Myofascial trigger point injection; dry needling; TP injection; trigger point block

\section*{Principle}
A trigger point is a hyperirritable taut band of skeletal muscle that generates local and referred pain. Injecting anesthetic into the point mechanically disrupts the taut band, releases muscle spasm, and blocks local nociceptive signaling, with corticosteroid added when inflammation is thought to contribute.

\section*{History}
Myofascial trigger points were described by Janet Travell in the mid-20th century, and injection-based treatment became a widely used technique for myofascial pain.

\section*{Why It Is Done}
Ordered for patients with myofascial pain syndrome, fibromyalgia-related focal pain, chronic neck and back pain, tension headache, and localized muscle spasm that has not responded to physical therapy, heat, massage, or analgesics.

\section*{Is This a Primary or Secondary Test or Procedure}
A secondary, adjunctive treatment for myofascial pain; it is therapeutic rather than diagnostic.

\section*{How It Is Performed}
The trigger point is identified by palpation of a taut band with a localized twitch response. The skin is cleaned, and a small-gauge needle is inserted into the trigger point, with the needle withdrawn slightly and redirected several times to elicit twitch responses before the anesthetic is injected.

\section*{Preparation}
Informed consent, review of anticoagulant use and bleeding risk, and identification of the target trigger point are needed. The patient should be positioned comfortably and should report if the injection causes distant or radiating pain.

\section*{Contraindications and Precautions}
Local skin infection, anticoagulation or bleeding disorder, and known allergy to the injected agents are contraindications. Precautions include avoiding direct injection into the lung, pleura, or major neurovascular structures and limiting corticosteroid use in patients with diabetes or immune compromise.

\section*{Risks}
Risks are generally low and include transient soreness, bleeding or bruising, infection, nerve injury, and pneumothorax when injecting near the chest wall.

\section*{Aftercare and Recovery}
The patient is observed briefly and advised that transient soreness is common. Physical therapy, stretching, and posture correction are encouraged to prevent recurrence, since injection alone provides temporary relief.

\section*{Outcome and Follow-up}
Immediate relief of the local and referred pain with improved muscle relaxation indicates a successful injection. Persistence or early recurrence of pain suggests the need for repeat injection combined with rehabilitative therapy or further evaluation.

\section*{Expected Results}
Complete or substantial relief of the injected trigger point and its referred pain, with normal muscle function, indicates a good response.

\section*{Follow-up}
Repeat injections are spaced to limit steroid exposure, and physical therapy with stretching exercises is continued. Imaging or electrodiagnostic studies may be considered if pain persists or if an alternative diagnosis is suspected.

\section*{References}
\begin{enumerate}
  \item MedlinePlus. Trigger point injection. U.S. National Library of Medicine; 2025.
  \item Current Medical Diagnosis and Treatment. McGraw-Hill; 2025.
\end{enumerate}

\clearpage
